% ======================================================================
%  PRÉAMBULE — Rapport de stage opérateur
%  (packages, mise en page, couleurs, commandes personnalisées)
% ======================================================================

% ---- Langue et encodage ---------------------------------------------
\usepackage[T1]{fontenc}
\usepackage[utf8]{inputenc}
\usepackage[french]{babel}      % typographie française
\usepackage{lmodern}
\usepackage{csquotes}

% ---- Mise en page ---------------------------------------------------
\usepackage[a4paper,margin=2.5cm,headheight=15pt]{geometry}
\usepackage{setspace}
\onehalfspacing                 % interligne 1.5 (consigne ECC)
\usepackage{microtype}

% ---- Graphisme ------------------------------------------------------
\usepackage{graphicx}
\usepackage{xcolor}
\usepackage{caption}            % pour \captionof hors flottant
\usepackage{tikz}
\usetikzlibrary{shapes.geometric,arrows.meta,positioning,calc,chains,fit,backgrounds}
\usepackage{pgfplots}
\pgfplotsset{compat=1.18}
\usepackage{tcolorbox}
\tcbuselibrary{skins,breakable}

% ---- Tableaux -------------------------------------------------------
\usepackage{booktabs}
\usepackage{longtable}
\usepackage{array}
\usepackage{multirow}

% ---- Listes ---------------------------------------------------------
\usepackage{enumitem}
\setlist{itemsep=2pt,topsep=3pt}

% ---- Titres et en-têtes --------------------------------------------
\usepackage{titlesec}
\usepackage{fancyhdr}
\usepackage[nottoc]{tocbibind}  % TdM/figures/tables dans la table des matières
\usepackage[titletoc]{appendix}

% ---- Liens et bibliographie ----------------------------------------
\usepackage[hidelinks]{hyperref}
\usepackage[backend=biber,style=numeric,sorting=none]{biblatex}
\addbibresource{references.bib}

% ======================================================================
%  COULEURS
% ======================================================================
\definecolor{centralegreen}{RGB}{0,124,102}   % vert Centrale (accent titres)
\definecolor{reminexorange}{RGB}{240,160,20}   % orange Reminex/Managem
\definecolor{darkslate}{RGB}{45,52,60}         % gris ardoise
\definecolor{softgray}{RGB}{120,128,138}

% ======================================================================
%  STYLE DES TITRES (chapitres = "PARTIE")
% ======================================================================
% Les chapitres numérotés sont intitulés « Partie »
\addto\captionsfrench{\renewcommand{\chaptername}{Partie}}

\titleformat{\chapter}[display]
  {\normalfont\huge\bfseries\color{darkslate}}
  {\filright\large\color{reminexorange}\MakeUppercase{\chaptertitlename}~\thechapter}
  {4pt}
  {\titlerule[1.1pt]\vspace{3pt}\filright}
  [\vspace{2pt}{\color{reminexorange}\titlerule[1.1pt]}]
\titlespacing*{\chapter}{0pt}{-18pt}{16pt}

\titleformat{\section}
  {\normalfont\large\bfseries\color{centralegreen}}{\thesection}{0.8em}{}
\titleformat{\subsection}
  {\normalfont\normalsize\bfseries\color{darkslate}}{\thesubsection}{0.7em}{}

% ======================================================================
%  EN-TÊTES ET PIEDS DE PAGE
% ======================================================================
\pagestyle{fancy}
\fancyhf{}
\renewcommand{\headrulewidth}{0.4pt}
\renewcommand{\footrulewidth}{0pt}
\fancyhead[L]{\small\color{softgray}Rapport de stage opérateur}
\fancyhead[R]{\small\color{softgray}Reminex --- Groupe Managem}
\fancyfoot[C]{\thepage}

\fancypagestyle{plain}{%
  \fancyhf{}\renewcommand{\headrulewidth}{0pt}\fancyfoot[C]{\thepage}}

% ======================================================================
%  COMMANDES PERSONNALISÉES
% ======================================================================

% --- Encadré "présenté par un ingénieur" (source orale) --------------
\newtcolorbox{presente}[1][]{%
  enhanced, breakable, colback=reminexorange!7, colframe=reminexorange!70,
  boxrule=0.4pt, arc=2pt, left=8pt, right=8pt, top=6pt, bottom=6pt,
  fonttitle=\bfseries, #1}

% --- Emplacement de photo (compile sans fichier externe) -------------
% Usage :  \photofig[largeur]{légende}{label}
% Pour insérer la vraie photo : remplacer le bloc \placeholderbox par
%   \includegraphics[width=<largeur>]{figures/mon_image.jpg}
\newcommand{\placeholderbox}[1]{%
  \fbox{\begin{minipage}[c][42mm][c]{#1}%
    \centering\color{softgray}\itshape\footnotesize
    [\,Emplacement réservé à une photo\,]\\[1mm]
    \tiny(remplacer par \textbackslash includegraphics)%
  \end{minipage}}}

\newcommand{\photofig}[3][0.68\linewidth]{%
  \begin{figure}[htbp]\centering
    \placeholderbox{#1}%
    \caption{#2}\label{#3}
  \end{figure}}

% --- Sigle mis en valeur ---------------------------------------------
\newcommand{\sig}[1]{\textsc{#1}}
